\vspace*{-0.6cm}
\cadlag es un sistema que calcula, para nueve ligas de fútbol y tres copas continentales, la probabilidad de que cada equipo termine en cada puesto de la tabla o llegue a cada ronda, y la publica en un sitio que se actualiza cada dos horas. El nombre viene de las trayectorias \emph{continuas por la derecha con límite por la izquierda} de los procesos estocásticos: la fuerza de un equipo cambia de a poco y salta cuando se juega un partido. Este documento explica la teoría, la implementación y los resultados, incluidos los negativos.

El núcleo es un \emph{modelo de espacio de estados}. Cada equipo tiene fuerzas que evolucionan como un paseo aleatorio dentro de la temporada y revierten a la media entre temporadas. Cada partido las observa a través de una verosimilitud de conteo. El modelo oficial separa crear ocasiones de convertirlas: cuatro fuerzas por equipo (crear y evitar tiros al arco, convertirlos y evitar que se los conviertan), con tiros al arco binomiales negativos y goles binomiales sobre los tiros. Por \emph{adelgazamiento}, la marginal de los goles es otra binomial negativa en forma cerrada. El estado se filtra partido a partido con una actualización de Laplace que resuelve toda la no linealidad en un problema de Newton de dimensión cuatro. Los hiperparámetros se ajustan por verosimilitud \emph{prequential} con gradientes exactos a través de todo el filtro, en JAX.

Sobre ese filtro se simulan 20.000 temporadas por liga con las reglas reales: desempates recursivos, partidos de definición, playoffs de ascenso, zonas y promedios del fútbol argentino. Cada temporada sortea primero las fuerzas y después los partidos. Esto es la predictiva posterior; fijar las fuerzas en su media subestima la varianza de los puntos a la mitad en pretemporada. Con las mismas simulaciones se mide la importancia de cada partido como información mutua, se calcula qué necesita cada equipo con un remuestreo condicional exacto y se estiman probabilidades de una en un millón con muestreo por importancia. El laboratorio del sitio condiciona en resultados elegidos por el usuario reponderando las 20.000 temporadas en el navegador, sin volver a simular.

Para las copas hace falta una escala común entre ligas. La capa continental parte de la fuerza doméstica de cada equipo, linealizada en goles, y le suma el nivel de su liga en el continente y un residuo propio. Se filtra con una aproximación de densidad supuesta factorizada que cuesta $O(1)$ por partido. La misma geometría aparece al comparar ligas con partidos de copa y ascensos: la varianza de la diferencia entre dos ligas es la \emph{resistencia efectiva} entre sus nodos en una red eléctrica.

Fuera de muestra, el modelo de dos latentes mejora la log-verosimilitud del marcador frente al modelo de goles clásico en el conjunto de nueve ligas ($+0{,}0036\pm0{,}0006$ nats por partido, $z=5{,}9$) y queda a $0{,}004$ de RPS de Bet365 en la Premier. La predictiva posterior lleva la cobertura de los intervalos de puntos de 76\,\% a 85\,\%, para un nominal de 90\,\%. Anclar los equipos a su liga mejora los partidos UEFA ($+0{,}030\pm0{,}004$, $z=7{,}2$) y los de Libertadores y Sudamericana ($+0{,}029\pm0{,}007$). Las ideas más atractivas a priori dan resultados negativos limpios: estilos latentes, intransitividad, contagio de goles y señales del mercado de pases.

\vspace{0.4em}
\textsf{\bfseries\small\color{cdpurple}Palabras clave:} modelo de espacio de estados; filtro de Laplace; Dixon–Coles; binomial negativa adelgazada; verosimilitud prequential; predictiva posterior; Monte Carlo; muestreo por importancia; información mutua; resistencia efectiva; densidad supuesta; reglas de puntuación propias; fútbol.
